% Part: propositional-logic
% Chapter: syntax-and-semantics
% Section: introduction

\documentclass[../../../include/open-logic-section]{subfiles}

\begin{document}

\olfileid{pl}{syn}{int}

\olsection{પરિચય}

વિધાનાત્મક તર્કશાસ્ત્રમાં $\lnot$, $\land$, $\lor$, $\lif$ અને
$\liff$ એ વિધાનાત્મક સંયોજકો વડે !!{propositional variable}sમાંથી
બનાવેલાં !!{formula}sનો અભ્યાસ થાય છે. સહજ રીતે કહીએ તો,
!!a{propositional variable}~$p$ સત્ય કે અસત્ય હોઈ શકે એવા વાક્ય
અથવા વિધાનને રજૂ કરે છે. !!a{formula}માંના !!{propositional variable}નું
``સત્યમૂલ્ય'' નક્કી થાય એટલે વિધાનાત્મક સંયોજકો વડે તેમાંથી બનાવેલાં
!!{formula}sનું સત્યમૂલ્ય પણ નક્કી થાય છે. આપણે
વિધાનાત્મક તર્કશાસ્ત્રને \emph{સત્યતાફલનલક્ષી} કહીએ છીએ, કારણ કે તેનો
અર્થવિચાર સત્યમૂલ્યોનાં વિધેયો વડે અપાય છે. ખાસ કરીને, વિધાનાત્મક
તર્કશાસ્ત્રમાં સત્ય અને અસત્યના બીજા કોઈ નિર્ધારણને ધ્યાનમાં લેતા નથી;
દાખલા તરીકે, કોઈ વાત માત્ર સંજોગવશાત્ સત્ય હોવાને બદલે અનિવાર્યપણે
સત્ય છે કે નહિ, તે સત્ય છે એવું જાણીતું છે કે નહિ, અથવા તે હમણાં
સત્ય છે કે પહેલાં સત્ય હતી કે પછી સત્ય થશે. આપણે માત્ર સત્ય
($\True$) અને અસત્ય ($\False$) એ બે સત્યમૂલ્યો લઈએ છીએ; એટલે કોઈ
કથન ન તો સત્ય ન તો અસત્ય, અથવા માત્ર અડધું સત્ય હોઈ શકે એવી શક્યતા
ચર્ચામાંથી બહાર રાખીએ છીએ. વળી, આપણે માત્ર એવા સંયોજકો પર ધ્યાન
આપીએ છીએ જેમનાથી બનેલા !!a{formula}નું સત્યમૂલ્ય તેના ભાગોનાં
સત્યમૂલ્યો વડે સંપૂર્ણપણે નક્કી થાય છે (અને, માનો કે, તેના અર્થ વડે
નહિ). ખાસ કરીને, અંગ્રેજી ભાષાનાં શરતી વાક્યોનું સત્યમૂલ્ય આ અર્થમાં
સત્યતાફલનલક્ષી છે કે નહિ તે વિવાદાસ્પદ છે. સત્યમૂલ્યલક્ષી પ્રેરણ
$\lif$ એવું છે; બીજા તર્કશાસ્ત્રો સત્યતાફલનલક્ષી ન હોય એવાં શરતી
વિધાનોનો અભ્યાસ કરે છે.

સત્યતાફલનલક્ષી વિધાનાત્મક તર્કશાસ્ત્રનો સિદ્ધાંત અને અધિસિદ્ધાંત
વિકસાવવા માટે પહેલાં તેના પદોની વાક્યરચના અને અર્થવિચાર વ્યાખ્યાયિત
કરવાં પડે. સંયોજકો વડે !!{propositional variable}sમાંથી
!!{formula}s રચવાની એક રીત આપણે વર્ણવીશું. બીજી વ્યાખ્યાઓ પણ શક્ય
છે. બીજી પદ્ધતિઓ જુદા સંકેતો પસંદ કરશે, સંયોજકોના જુદા ગણને મૂળભૂત
ગણશે અને કૌંસ જુદી રીતે વાપરશે (અથવા કહેવાતી પોલિશ સંજ્ઞાપદ્ધતિની
જેમ જરા પણ નહિ વાપરે). છતાં, બધી રીતોમાં એક વાત સામાન્ય છે: રચનાના
નિયમો !!{formula}sના ગણને \emph{અનુમાનાત્મક રીતે} વ્યાખ્યાયિત કરે છે.
નિયમો યોગ્ય રીતે આપ્યા હોય, તો દરેક પદ રચનાના નિયમો અનુસાર મૂળભૂત રીતે
એક જ રીતે બની શકે. પદોને \emph{એકમાત્ર વાંચી શકાય} એવી અનુમાનાત્મક
વ્યાખ્યા મળવાથી એ જ રીત---અનુમાનાત્મક વ્યાખ્યા---વાપરીને આપણે આ
પદોને અર્થ આપી શકીએ છીએ.

પદોને અર્થ આપવાનું ક્ષેત્ર અર્થવિચાર છે. વિધાનાત્મક તર્કશાસ્ત્રના
અર્થવિચારમાં !!a{valuation}માં થતો સંતોષ મુખ્ય ખ્યાલ છે.
!!^a{valuation}~$\pAssign{v}$ !!{propositional variable}sને સત્યમૂલ્યો
$\True$, $\False$ નિયુક્ત કરે છે. દરેક !!{valuation}, દરેક
!!{formula}~$!A$ માટે સત્યમૂલ્ય $\pValue{v}(!A)$ નક્કી કરે છે.
$\pValue{v}(!A)=\True$ હોય તો અને તો જ !!^a{formula},
!!a{valuation}~$\pAssign{v}$માં સંતોષાય છે---આને આપણે
$\pSat{v}{!A}$ તરીકે લખીએ છીએ. આ સંબંધને પણ $!A$ની રચના પર અનુમાન
વડે વ્યાખ્યાયિત કરી શકાય છે: તાર્કિક સંયોજકોનાં સત્યવિધેયો વાપરીને,
માનો કે, $!A\land !B$નો સંતોષ $!A$ અને $!B$ના સંતોષ (અથવા
અસંતોષ)ના આધારે વ્યાખ્યાયિત થાય છે.

વાક્યો માટેના સંતોષસંબંધ $\pSat{v}{!A}$ને આધારે પછી આપણે પુનરુક્તિ,
ફલિતતા અને સંતોષ્યતાના મૂળ અર્થવિચારી ખ્યાલો વ્યાખ્યાયિત કરી શકીએ
છીએ. દરેક !!{valuation} તેને સંતોષે, એટલે કે દરેક $\pAssign{v}$ માટે
$\pValue{v}(!A)=\True$ હોય, તો !!^a{formula} પુનરુક્તિ છે અને આપણે
$\Entails !A$ લખીએ છીએ. જો $\Gamma$નાં બધાં !!{formula}sને સંતોષતું
દરેક !!{valuation}, $!A$ને પણ સંતોષે, તો !!{formula}sનો ગણ તેને
ફલિત કરે છે અને આપણે $\Gamma\Entails !A$ લખીએ છીએ. અને કોઈ એક
!!{valuation}, !!{formula}sના ગણમાંનાં બધાં !!{formula}sને એકસાથે સંતોષે તો
તે ગણ સંતોષ્ય છે. !!{formula}s અનુમાનાત્મક રીતે વ્યાખ્યાયિત છે અને
સંતોષ પણ !!{formula}sની રચના પર અનુમાન વડે વ્યાખ્યાયિત થાય છે; તેથી
અર્થવિચારના ગુણધર્મો સાબિત કરવા અને વ્યાખ્યાયિત અર્થવિચારી ખ્યાલોને
પરસ્પર જોડવા આપણે અનુમાન વાપરી શકીએ છીએ.


\end{document}
